
We state three results: the cursors of a run do not depend on its randomness
(Lemma~\ref{lem:cursors}), replay reproduces a recorded run exactly
(Theorem~\ref{thm:replay}), and resuming after a crash continues the run with the
correct distribution (Corollary~\ref{cor:resume}).
Each comes in a pointwise form, quantified over every tree and every tape, which is
mechanized in Lean~4 (Section~\ref{sec:mechanization}), and a probabilistic form, which
follows on paper from the pointwise one.
Throughout we assume:
\begin{description}[noitemsep]
  \item[(A1)] $\mathit{Val}$ is countable and each $\meas{p}$ is a probability measure on it.
  \item[(A2)] Replay and resume run the same program, from the same starting cursor, as the
    recorded run.
  \item[(A3)] The journal is not corrupted or reordered; a crash keeps a prefix of it.
\end{description}

\begin{lemma}[Cursors]
\label{lem:cursors}
For every tree $T$, tape $\omega$, and cursor $c$, the $i$-th entry of the journal of
$\mathsf{rec}(T, \omega, c)$ has cursor $c+i-1$, whatever the responses (only the number of
entries depends on them).
\end{lemma}

\begin{proof}
By induction on $T$: each call appends one entry at the current cursor and advances it by one.
\end{proof}

\begin{theorem}[Replay soundness]
\label{thm:replay}
For every tree $T$, tape $\omega$, and cursor $c$, let $(t, J) = \mathsf{rec}(T, \omega, c)$.
Then
\[
  \mathsf{rep}(T, J) \;=\; (t, J).
\]
Consequently:
\begin{enumerate}[label=(\roman*), noitemsep]
  \item \textbf{(Trace equivalence)} replaying the journal reproduces the recorded trace,
    including every cursor, prompt, and response, and the journal itself;
  \item \textbf{(Determinism)} replay is a function of the journal alone and reads no
    randomness, so given $J$ it produces $t$ with probability~1;
  \item \textbf{(Equiprobability)} for every set $\Phi$ of traces,
    \[
      \mathbb{P}\bigl[\,\mathsf{rec}(T,\omega,c) \text{ has trace in } \Phi\,\bigr]
      \;=\;
      \mathbb{P}\bigl[\,\mathsf{rep}(T, J_\omega) \text{ has trace in } \Phi\,\bigr],
    \]
    where $J_\omega$ is the journal of $\mathsf{rec}(T,\omega,c)$; that is, recording and then
    replaying induces the same distribution over traces as recording alone, and by
    Lemma~\ref{lem:tape} a trace $o(c,p_1,v_1)\cdots o(c{+}n{-}1,p_n,v_n)\,\mathsf{done}(v)$
    of $T$ has probability $\prod_i \meas{p_i}(\{v_i\})$ under both.
\end{enumerate}
\end{theorem}

\begin{proof}
The identity is proved by induction on $T$.
For $\mathsf{halt}(v)$ both sides are $(\mathsf{done}(v), \varepsilon)$.
For $\mathsf{draw}(p, k)$, recording at cursor $c$ produces the entry $(c, p, v)$ with
$v = \omega(c,p)$ and continues as $\mathsf{rec}(k(v), \omega, c+1) = (t', J')$, so
$(t, J) = (o(c,p,v)\,t',\; (c,p,v)\,J')$.
Replay reads the entry $(c, p, v)$, emits $o(c, p, v)$, and continues as
$\mathsf{rep}(k(v), J')$, which equals $(t', J')$ by the induction hypothesis for the subtree
$k(v)$.
Part~(i) is the first component of the identity, and part~(ii) holds because $\mathsf{rep}$
does not take a tape.
For part~(iii), the identity holds for every $\omega$, so the two events are the same subset
of $\Omega$ and have the same probability; the product formula is Lemma~\ref{lem:tape}.
\end{proof}

\noindent
Part~(iii) is immediate from part~(i): it holds for every law on tapes, not only $\mathbb{P}$,
and the same is true of deterministic replay under any input distribution.
The probabilistic content of the paper lies elsewhere: in Lemma~\ref{lem:tape}, which says that
the tape semantics realizes sampling from the oracle, and in Corollary~\ref{cor:resume}(ii),
which says that recovery with fresh randomness preserves the distribution of a run.

\begin{figure}[t]
  \centering
  \fitwidth{\begin{tikzpicture}[
  >=Stealth,
  every node/.style={align=center, font=\small},
  rec/.style={draw, circle, inner sep=0pt, minimum size=17pt, font=\scriptsize},
  rep/.style={draw, circle, densely dashed, fill=black!10, inner sep=0pt, minimum size=17pt, font=\scriptsize},
  jb/.style={draw, rounded corners=1pt, minimum width=1.3cm, minimum height=0.46cm, inner sep=1pt, font=\scriptsize},
  jv/.style={->, densely dashed, thin, gray},
  lbl/.style={font=\scriptsize\itshape},
]

\foreach \i/\x in {1/2, 2/3.5, 3/5, 4/6.5} { \node[rec] (s\i) at (\x,0) {$s_{\i}$}; }
\draw[->] (s1)--(s2); \draw[->] (s2)--(s3); \draw[->] (s3)--(s4);

\draw[->] (s4) -- (7.4,0);
\draw[very thick] (7.45,-0.22)--(7.95,0.22);
\draw[very thick] (7.95,-0.22)--(7.45,0.22);
\node[font=\footnotesize\bfseries] at (7.7,0.5) {crash};
\node[lbl, anchor=west] at (8.15,0) {in-memory state lost};

\draw[thin] (1.4,0.55)--(7.0,0.55);
\draw[thin] (1.4,0.42)--(1.4,0.6);
\draw[thin] (7.0,0.42)--(7.0,0.6);
\node[lbl] at (4.2,0.85) {original recorded run (RECORD)};

\foreach \i/\x in {1/2, 2/3.5, 3/5, 4/6.5} {
  \node[jb] (j\i) at (\x,-1.5) {$(k_{\i},v_{\i})$};
  \draw[jv] (s\i.south) -- (j\i.north);
}
\node[lbl, anchor=west] at (7.3,-1.5) {$J^*$: one entry appended per step (durable)};

\foreach \i/\x in {1/2, 2/3.5, 3/5, 4/6.5} {
  \node[rep] (r\i) at (\x,-3.5) {$s_{\i}$};
  \draw[jv] (j\i.south) -- (r\i.north);
}
\draw[->, densely dashed] (r1)--(r2);
\draw[->, densely dashed] (r2)--(r3);
\draw[->, densely dashed] (r3)--(r4);

\draw[densely dotted, thick] (7.3,-3.95)--(7.3,-3.05);
\draw[->] (r4) -- (7.25,-3.5);
\node[rec] (s5) at (8.3,-3.5) {$s_5$};
\draw[->] (7.35,-3.5) -- (s5);
\node[lbl] at (7.3,-2.75) {$[O\text{-}resume]$};

\draw[thin] (1.4,-4.05)--(7.0,-4.05);
\draw[thin] (1.4,-3.9)--(1.4,-4.08);
\draw[thin] (7.0,-3.9)--(7.0,-4.08);
\node[lbl, align=center] at (4.2,-4.45) {replay prefix: reproduces the original\\(Theorem~\ref{thm:replay})};
\draw[thin] (7.8,-4.05)--(8.8,-4.05);
\draw[thin] (7.8,-3.9)--(7.8,-4.08);
\draw[thin] (8.8,-3.9)--(8.8,-4.08);
\node[lbl, align=center] at (8.3,-4.45) {new\\(record)};

\draw[->, thick] (1.0,-0.35) -- (1.0,-3.15);
\node[lbl, rotate=90, anchor=south] at (0.72,-1.75) {resume: replay $J^*$ from the start};
\end{tikzpicture}
}
  \caption{Crash recovery (Corollary~\ref{cor:resume}).
    Calls $s_1 \ldots s_4$ run in record mode and append their outcomes to the journal; a crash
    at $s_5$ loses all in-memory state but keeps the journal prefix. Resume replays the prefix
    with probability~1 (\textsc{[O-rep]}), then switches to record mode
    (\textsc{[O-resume]}) for the rest of the run.}
  \label{fig:crash-recovery}
\end{figure}

\subsection{Crash Recovery}

A crash loses the agent's in-memory state but, by (A3), keeps a prefix of the journal.
Recovery resumes the program on that prefix with whatever randomness is available afterwards.
To state the result, define the \emph{splice} of two tapes at cursor $m$ by
\[
  \mathsf{splice}_m(\omega, \omega')(c, p) \;=\;
  \begin{cases}
    \omega(c, p) & c < m,\\
    \omega'(c, p) & c \geq m.
  \end{cases}
\]

\begin{corollary}[Crash recovery]
\label{cor:resume}
For every tree $T$, tapes $\omega$ and $\omega'$, cursor $c$, and $n \in \mathbb{N}$, let $J$
be the journal of $\mathsf{rec}(T, \omega, c)$ and $J_{<n}$ its first $n$ entries.
Then
\[
  \mathsf{res}(T,\; J_{<n},\; \omega',\; c) \;=\; \mathsf{rec}\bigl(T,\; \mathsf{splice}_{c+n}(\omega, \omega'),\; c\bigr).
\]
Consequently:
\begin{enumerate}[label=(\roman*), noitemsep]
  \item \textbf{(Same randomness)} resuming with $\omega' = \omega$ reproduces the uninterrupted
    run exactly, wherever the crash occurred;
  \item \textbf{(Fresh randomness)} if $\omega'$ is drawn from $\mathbb{P}$ independently of
    $\omega$, the resumed run has the same distribution as an uninterrupted record run, and
    conditioned on the surviving prefix $J_{<n}$, its continuation has the same distribution
    as the original run's continuation conditioned on the same prefix.
\end{enumerate}
\end{corollary}

\begin{proof}
The identity is proved by induction on $T$, generalizing over $c$ and $n$, using one auxiliary
fact: a record run from cursor $c$ reads the tape only at cursors $\geq c$, so two tapes that
agree from $c$ on give the same run.
If $n = 0$, resuming on the empty journal is recording on $\omega'$, and
$\mathsf{splice}_c(\omega, \omega')$ agrees with $\omega'$ from $c$ on.
If $n > 0$ and $T = \mathsf{draw}(p, k)$, the first surviving entry is $(c, p, v)$ with
$v = \omega(c, p) = \mathsf{splice}_{c+n}(\omega,\omega')(c, p)$, so both sides make the same
first call and continue with $k(v)$ from cursor $c+1$ with $n-1$ surviving entries, where the
induction hypothesis applies.

For~(i), $\mathsf{splice}_m(\omega, \omega) = \omega$.
For~(ii), the coordinates of $\mathsf{splice}_{c+n}(\omega, \omega')$ are independent with
the right marginals, since those at cursors below $c+n$ come from $\omega$, the rest come from
$\omega'$, and $\omega$ and $\omega'$ are independent; so the spliced tape has law $\mathbb{P}$
and the resumed run has the record distribution.
By Lemma~\ref{lem:cursors} the surviving prefix reads only coordinates at cursors below $c+n$,
on which the two tapes agree, so it is the same random variable under $\omega$ and under the
spliced tape.
The pair (prefix, whole run) therefore has the same joint law under both tapes, and so does the
conditional law of the run given the prefix.
\end{proof}

\begin{remark}[When the crash point is random]
Corollary~\ref{cor:resume} fixes the crash point $n$.
The identity holds for every $n$, so it also holds when $n$ is random.
Part~(ii) then still holds if the crash point is a \emph{stopping time}, meaning whether the
crash happens after $k$ entries depends only on the coordinates at cursors below $c+k$, and is
independent of $\omega'$: replacing the coordinates after a stopping time by independent fresh
ones does not change the law of a tape whose cursor rows are independent and identically
distributed, which is the strong Markov property.
It fails if the crash depends on a response that was drawn but not yet journaled.
For example, let $T$ make one call whose response is a fair coin and return it, and let the
process crash before journaling exactly when the response is $0$.
Then the resumed run returns $1$ with probability $3/4$, not $1/2$.
A recovery guarantee in distribution therefore needs the response to be journaled before
anything that depends on it can crash the agent, which is the order \textsc{[O-rec]} enforces.
\end{remark}

\noindent
For a single agent every prefix of its journal is a valid recovery point, so no
consistent-cut condition is needed; Section~\ref{sec:limitations} explains why concurrent
agents need one.

\subsection{Mechanization}
\label{sec:mechanization}

Lemma~\ref{lem:cursors}, the identity of Theorem~\ref{thm:replay} with part~(i) and the
pointwise core of part~(iii), and the identity of Corollary~\ref{cor:resume} with part~(i) are
mechanized in Lean~4; part~(ii) of the theorem holds by the type of the replay function, which
takes no tape (\texttt{lean/AgenticalT1.lean}, checked with Lean v4.30.0).
The development represents values, prompts, and cursors by natural numbers, programs by the
trees of Section~\ref{sec:syntax}, and tapes by functions
$\mathbb{N} \to \mathbb{N} \to \mathbb{N}$, and quantifies over every tree, tape, starting
cursor, and crash point.
It contains no unproven steps, and \texttt{\#print axioms} reports only Lean's standard axioms
(\texttt{propext}, and \texttt{Quot.sound} for the crash-recovery theorems).
The translation from terms to trees and the measure-theoretic layer, Lemma~\ref{lem:tape} and
the distributional parts of Theorem~\ref{thm:replay}(iii) and Corollary~\ref{cor:resume}(ii),
are on paper; their proofs use the pointwise identities, Lemma~\ref{lem:cursors}, and the
independence and marginals of the tape coordinates.
